% part: intuitionistic-logic
% chapter: propositions-as-types
% section: types

\documentclass[../../../include/open-logic-section]{subfiles}

\begin{document}

\olfileid[hi]{int}{pty}{typ}

\olsection{प्रकार}

$(MN)$ जैसा प्रमाण पद अपने आप उस प्रमाण को अद्वितीय रूप से निर्धारित नहीं
करता जिसका वह प्रमाण पद है। कारण यह है कि मुक्त चर, अर्थात् अभिगृहीतों के
प्रमाण पद, संगत अभिगृहीतों के !!{formula}s को विनिर्दिष्ट नहीं करते। किंतु यदि
हम इन्हें किसी प्रसंग~$\Gamma$ में विनिर्दिष्ट करें, तो प्रमाण अद्वितीय रूप से
निर्धारित \emph{होता है}।

\begin{defn}
किसी प्रमाण पद~$N$ का \emph{प्रसंग} ऐसा समुच्चय $\Gamma$ है जिसमें~$N$ के
प्रत्येक मुक्त चर के लिए ठीक एक युग्म $x : !A$ हो।
\end{defn}

ध्यान दें कि परिभाषा यह आवश्यक नहीं करती कि $\Gamma$ में $x:!A$ \emph{केवल}
~$N$ के मुक्त चरों के लिए हो। इस प्रकार, उदाहरणतः $\Gamma = \{x:!A, y:!B,
z:!C\}$, $\pair{x}{y}$ का प्रसंग है।

\begin{defn}\ollabel{defn:type}
  मान लें $\Gamma$,~$N$ का प्रसंग है।~$N$ का \emph{प्रकार} (~$\Gamma$ के
  सापेक्ष) आगमनिक रूप से इस प्रकार परिभाषित है:
  \begin{enumerate}
  \item $x$ का प्रकार $!A$ तभी और केवल तभी है जब $x:!A \in \Gamma$।
  \item यदि $x:!A, \Gamma$ के सापेक्ष $N$ का प्रकार~$!B$ है, तो
  $\lambd[\typeof{x}{!A}][N]$ का प्रकार $!A \lif !B$ है।
  \item यदि $N$ का प्रकार $!A \lif !B$ और $M$ का प्रकार~$!A$ है, तो $(NM)$
  का प्रकार~$!B$ है।
  \item यदि $N$ का प्रकार $!A$ और $M$ का प्रकार~$!B$ है, तो $\pair{N}{M}$ का
  प्रकार $!A \land !B$ है।
  \item यदि $N$ का प्रकार~$!A_1 \land !A_2$ है, तो $\proj{i}{N}$ का
  प्रकार~$!A_i$ है।
  \item यदि $N$ का प्रकार $!A$ है, तो $i=1$ होने पर $\inj[!B]{i}{!M}$ का
  प्रकार $!A \lor !B$ और $i=2$ होने पर $!B \lor !A$ है।
  \item यदि $M$ का प्रकार $!A \lor !B$, $x:!A, \Gamma$ के सापेक्ष $N_1$ का
  प्रकार $!C$, और $y:!B, \Gamma$ के सापेक्ष $N_2$ का प्रकार~$!C$ है, तो
  $\dcase{M}{x_1}{N_1}{x_2}{N_2}$ का प्रकार~$!C$ है।
  \item यदि $N$ का प्रकार~$\lfalse$ है, तो $\abort{!A}{N}$ का प्रकार~$!A$ है।
  \end{enumerate}
\end{defn}

\begin{prop}
प्रत्येक प्रमाण पद~$N$ का उसके प्रसंग~$\Gamma$ के सापेक्ष ठीक एक प्रकार है।
\end{prop}

\begin{proof}
  ~ $N$ पर आगमन से।
\end{proof}

यदि $\Gamma$ के सापेक्ष $N$ का प्रकार~$!A$ है, तो हम $\Gamma \Sequent N :
!A$ लिखते हैं। \olref{defn:type} के उपवाक्य परिचित लगने चाहिए: वे एक छोटे अंतर
के साथ ठीक \Log{tN2i} के नियम हैं—प्रसंग~$\Gamma$ पूरे समय समान है। परिभाषा को
\olref{tab:tN2ip} के नियमों वाले कलन \Log{tN3i} के रूप में संहिताबद्ध कर सकते
हैं।

\subfile{rules-tN3}

प्रमाण पद $\lambd$-कलन के एक रूप—\emph{गुणनफलों, योगों और रिक्त प्रकार वाला
प्रकारित $\lambd$-कलन}—के पद हैं। चिरसम्मत $\lambd$-कलन में केवल चरों से बने
पद, \emph{$\lambd$-अमूर्तन}~$\lambd[x][N]$ और अनुप्रयोग $(MN)$ हैं; उसमें
प्रकारित $\lambd$-अमूर्तन $\lambd[\typeof{x}{!A}][N]$ का~$!A$ द्वारा अंकन नहीं
है। $\lambd$-कलन अभिकलन का ट्यूरिंग-पूर्ण निदर्श है (अर्थात् उसमें प्रत्येक
आंशिक अभिकलनीय फलन को किसी पद से निरूपित किया जा सकता है)। प्रमाण पदों और
प्रकारण नियमों द्वारा दिया गया $\lambd$-कलन का रूप अभिकलन का प्रतिबंधित निदर्श
है, किंतु मूल विचार वही हैं। प्रकार-निर्धारण उसे प्रतिबंधित करता है: केवल
\emph{सुप्रकारित} पद, अर्थात् किसी प्रसंग~$\Gamma$ के सापेक्ष प्रकार वाले प्रमाण
पद, वैध हैं। उदाहरणतः $\lambd[x][(xx)]$ चिरसम्मत अप्रकारित $\lambd$-कलन में
वैध पद है, किंतु $\lambd[\typeof{x}{!A}][(xx)]$ किसी भी $!A$ के लिए सुप्रकारित
नहीं है। कारण यह है कि $(xx)$ का प्रकार होने के लिए प्रकारण नियमों के अनुसार
पहले~$x$ का प्रकार $!A \lif !B$ और दूसरे $x$ का प्रकार~$!A$ होना चाहिए; यह
असंभव है क्योंकि $!A$,~$!A \lif !B$ का उचित उप-!!{formula} है।

सहजतः प्रकारों को उस प्रकार के किसी पद द्वारा नामित वस्तु का प्रकार समझ सकते
हैं। अतः $!A \land !B$ प्रकार का पद ऐसा युग्म चुनता है जिसका पहला घटक~$!A$
और दूसरा घटक~$!B$ प्रकार का है, अर्थात् गुणनफल~$!A \times !B$ का
!!a{element}। $!A \lif !B$ प्रकार का पद, $!A$ प्रकार की वस्तुओं से~$!B$ प्रकार
की वस्तुओं में फलन है। $!A \lor !B$ प्रकार का पद या तो $!A$ प्रकार की कोई
वस्तु है या~$!B$ प्रकार की, और साथ में यह सूचना है कि कौन-सी। यह दो समुच्चयों
$!A$ और $!B$ के असंयुक्त सम्मिलन जैसा है, जिसे $\{\tuple{i,x} : i = 1
\text{ अथवा } 2, x \in !A \text{ यदि } i = 1, x \in !B \text{ यदि } i=2\}$
से परिभाषित किया गया है। प्रकार $\lfalse$ रिक्त समुच्चय है; अर्थात् किसी पद
का यह प्रकार हो तो वह कुछ नामित नहीं कर सकता। चूँकि प्राकृतिक निगमन
(\Log{tN3i} सहित) संगत है, खुली मान्यताओं के बिना~$\lfalse$ का कोई !!{proof}
नहीं है; अतः~$\lfalse$ प्रकार का कोई बंद प्रमाण पद (मुक्त चरों से रहित पद)
नहीं है।

प्रकारित $\lambd$-कलन के संकारक ऐसे पद उत्पन्न करते हैं जो सहजतः सही प्रकार
की वस्तु को नामित करते हैं, बशर्ते जिन पदों पर उन्हें लगाया गया है वे कुछ
निश्चित वस्तुओं को नामित करें। उदाहरणतः ``यदि $N_1 : !A$ और $N_2 : !B$, तो
$\pair{N_1}{N_2} : !A \land !B$'' कहता है कि यदि $N_1$,~$!A$ प्रकार की और
$N_2$,~$!B$ प्रकार की किसी वस्तु को नामित करता है, तो $\pair{N_1}{N_2}$,
$!A \land !B$ प्रकार की किसी वस्तु को नामित करता है।

\end{document}
